\begin{resumo}

Modelos epidemiológicos como o SIR oferecem uma leitura interpretável da difusão de conteúdos culturais em plataformas de streaming, mas operam sobre curvas de popularidade agregadas e não capturam a estrutura relacional entre músicas, artistas e gêneros que também condiciona essa difusão. Este trabalho investiga em que medida uma rede neural em grafo heterogêneo com componente temporal explícito melhora a modelagem da difusão musical em relação a esse baseline populacional. A partir do ecossistema de dados do MGD+ combinado a rankings diários do Spotify (Top 200 e Viral 50) no mercado brasileiro entre 2017 e 2021, constrói-se um grafo heterogêneo com 6.526 músicas, 1.701 artistas e 530 gêneros, conectados por relações de autoria, associação a gênero, coocorrência de gêneros e cotrajetória de popularidade. Sobre esse grafo treina-se o \textsc{MusicDiffusionGNN}, que combina um codificador HeteroGraphSAGE por \emph{snapshots} semanais com uma cabeça recorrente GRU, comparado ao baseline SIR sob dois regimes: reconstrução retroativa da curva completa (Modo~1) e previsão causal a $k$ semanas à frente usando apenas informação passada (Modo~2). O baseline SIR foi antes reproduzido com fidelidade ao trabalho de referência (RMSE de viralidade $0{,}0289$ e de sucesso $0{,}0471$, $1.981$ músicas), e o grafo construído apresentou estrutura topológica não aleatória, com \textit{clustering} elevado e comunidades de gênero coerentes. No Modo~1, o SIR obtém menor RMSE por ajustar dois parâmetros livres à curva completa de cada música, uma vantagem estrutural do regime \emph{in-sample}, embora a diferença se estreite nas semanas em que a música está de fato no chart. No Modo~2, que corresponde ao uso pretendido de um preditor, a GNN supera o SIR nos seis cenários avaliados (dois charts $\times$ três horizontes), com significância estatística forte (Wilcoxon pareado, $p$ entre $6{,}5\times10^{-4}$ e $5{,}4\times10^{-32}$) e menor erro em 64\% a 77\% das músicas, superando também o baseline de persistência nos horizontes mais longos. Esses resultados indicam que a GNN supera ambos os baselines nos horizontes mais longos, com vantagem sobre o SIR que cresce com o horizonte; se esse ganho decorre especificamente da estrutura relacional do grafo, e não apenas da maior flexibilidade de um preditor neural treinado frente a um ajuste de dois parâmetros, é uma questão que permanece em aberto. Como trabalho em andamento, permanecem a análise de interpretabilidade da GNN (ablação por tipo de aresta, importância de atributos e extração de análogos de $\beta$, $\gamma$ e $R_0$) e um baseline neural sem grafo, a serem incorporados na versão final da dissertação.

\end{resumo}
